\chapter{Конструкторская часть}

\section{Архитектура метода}

Разрабатываемый метод обеспечения дифференциальной приватности при генерации синтетических данных основан на Байесовских сетях и состоит из трех основных этапов:
\begin{enumerate}[label*=---]
	\item предварительная обработка входных данных;
	\item построение структуры Байесовской сети с дифференциальной приватностью;
	\item генерация синтетических данных на основе построенной сети.
\end{enumerate}

Архитектура системы представляет собой модульную структуру, где каждый компонент выполняет специфическую функцию в процессе генерации приватных синтетических данных.

На рисунке~\ref{img:idef1} представлена функциональная модель разрабытываемого метода.

\includeimage
{idef1} % Имя файла без расширения (файл должен быть расположен в директории inc/img/)
{f} % Обтекание (с обтеканием)
{h!} % Положение рисунка (см. wrapfigure из пакета wrapfig)
{1\textwidth} % Ширина рисунка
{Функциональная модель разрабатываемого метода} % Подпись рисунка

\section{Предварительная обработка данных}

Этап предварительной обработки является критически важным для качественной генерации синтетических данных. Система автоматически определяет типы данных и их характеристики, что обеспечивает корректную работу последующих алгоритмов.

Тип данных является важным компонентом домена атрибутов.
Метод поддерживает пять типов данных:
\begin{enumerate}[label*=---]
	\item Integer --- целые числа;
	\item Float --- числа с плавающей точкой;
	\item String --- строковые значения;
	\item DateTime --- временные метки.
\end{enumerate}

Система позволяет пользователям явно указывать типы данных. Если тип данных атрибута не указан пользователем, то метод выполняет автоматическое определение типа на основе анализа содержимого столбцов.

Алгоритм определения типа данных:
\begin{enumerate}
	\item Проверка на числовой тип (Integer/Float);
	\item Для нечисловых данных --- проверка на DateTime;
	\item Проверка на специальные форматы (SSN);
	\item По умолчанию --- String тип.
\end{enumerate}

Тип данных атрибута ограничивает его домен, который может быть дополнительно ограничен, если атрибут категоричен.
Например, в наборе данных студентов атрибут степени может принимать только значения $\{BS, BA, MS, PhD\}$.
Атрибут считается категориальным, если во входном наборе данных наблюдается ограниченное количество различных значений для этого атрибута (по умолчанию не более 20 уникальных значений).

\newpage
Математически категориальный признак определяется как:
\begin{equation}
    \text{is\_categorical}(A) = \begin{cases}
        \text{True}, & \text{если } |unique(A)| \leq threshold, \\
        \text{False}, & \text{иначе},
    \end{cases}
\end{equation}
где:
\begin{enumerate}[label*=---]
	\item $unique(A)$ --- множество уникальных значений атрибута $A$;
	\item $threshold$ --- пороговое значение.
\end{enumerate}


Этап предварительной обработки данных продемонстрирован на рисунке~\ref{img:anal-algos}:

\includeimage
{anal-algos} % Имя файла без расширения (файл должен быть расположен в директории inc/img/)
{f} % Обтекание (с обтеканием)
{h!} % Положение рисунка (см. wrapfigure из пакета wrapfig)
{0.9\textwidth} % Ширина рисунки
{Этап предварительной обработки данных} % Подпись рисунка
\newpage
\section{Построение структуры БС с ДП}

Построение структуры Байесовской сети с дифференциальной приватностью является ключевым этапом разработанного метода. Алгоритм использует жадный подход для поиска оптимальной структуры сети.

Алгоритм строит БС $\mathcal{N}$ из входного набора данных $D$, атрибутов $A$ и максимального количества родителей узла БС $k$. 

В алгоритме:
\begin{itemize}[label*=---]
	\item $V$ --- множество посещенных атрибутов;
	\item $\Pi$ --- подмножество $V$, которое станет родителями узла $X$;
	\item $\Omega$ --- множество всех возможных пар (атрибут, родители).
\end{itemize}

Атрибуты $\Pi$ определяются путем поиска максимального $I(X, \Pi)$, который вычисляется с добавлением шума Лапласа по формуле~(\ref{eq:dp}):
\begin{equation}\label{eq:dp}
	I(X_i; \Pi) = \sum_{x_i, \pi} P(x_i, \pi) \ln \frac{P(x_i, \pi)}{P(x_i) P(\pi)} + \text{Lap}\left(\frac{\Delta I}{n\varepsilon}\right),
\end{equation}
где:
\begin{enumerate}[label*=---]
	\item $\varepsilon$ --- параметр приватности;
	\item  $n$ --- количество записей входных данных;
	\item  $\Delta I$ --- чувствительность взаимной информации.
\end{enumerate}
Чувствительность взаимной информации вычисляется по формуле:
\begin{equation}
    \Delta I = \frac{2 \ln(n+1)}{n},
\end{equation}

%Процесс построения сети:
%%\item Инициализация пустой сети $\mathcal{N} = \emptyset$;
	%\item Случайный выбор корневого атрибута;
	%\item Итеративное добавление атрибутов с оптимальными родителями;
	%\item Применение экспоненциального механизма для приватного выбора.
%\end{enumerate}

%\begin{algorithm}
	%\label{alg:preproceses}
	%\caption{Построение БС с ДП}
	%\small
	%\SetKwInOut{Input}{Вход}
	%\SetKwInOut{Output}{Выход}
	
	%\Input{
	%	\Indp
	%	$D$~: датасет;\\
	%	$A$~: набор атрибутов;\\
	%	$k$~: максимальное количество родительских узлов;\\
	%	$\varepsilon$~: параметр приватности;\\
	%}
	%\Output{
	%	\Indp
	%	$\mathcal{N}$~: структура Байесовской сети.\\
	%}
	
	%\Begin{
	%	$\mathcal{N} \leftarrow \emptyset$ и $V \leftarrow \emptyset$ \\
	%	$X_1 \leftarrow$ случайный выбор из $A$ \\
	%	$\mathcal{N} \leftarrow \mathcal{N} \cup \{(X_1, \emptyset)\}$ и $V \leftarrow V \cup \{X_1\}$ \\
		
	%	\For{$i \gets 2 \text{ to } |A|$}{
		%	$\Omega \gets \emptyset$ \\
		%	$p \gets \min(k, |V|)$ \\
		%	\For{каждый $X \in A \setminus V$ и каждый $\Pi \subseteq V, |\Pi| \leq p$}{
			%	$\Omega \leftarrow \Omega \cup \{(X, \Pi)\}$ \\
			%	Вычислить $I(X; \Pi)$ с шумом Лапласа
		%	}
		%	$(X_i, \Pi_i) = \underset{(X,\Pi) \in \Omega}{\mathrm{argmax}} \tilde{I}(X; \Pi)$ \\
		%	$\mathcal{N} \leftarrow \mathcal{N} \cup \{(X_i, \Pi_i)\}$ \\
		%	$V \leftarrow V \cup \{X_i\}$
	%	}

	%	\Return $\mathcal{N}$
%	}
% \end{algorithm}
\newpage
Схема алгоритма для построения структуры БС с ДП представлена на рисунке~\ref{img:bs-algos}
\includeimage
{bs-algos} % Имя файла без расширения (файл должен быть расположен в директории inc/img/)
{f} % Обтекание (с обтеканием)
{h!} % Положение рисунка (см. wrapfigure из пакета wrapfig)
{\textwidth} % Ширина рисунка
{Схема алгоритма для построения структуры БС с ДП} % Подпись рисунка

% \subsubsection{Экспоненциальный механизм для приватного выбора}

Для обеспечения дифференциальной приватности при выборе оптимальной структуры сети используется экспоненциальный механизм. Вероятность выбора пары $(X, \Pi)$ определяется по формуле~(\ref{eq:c}):

\begin{equation}\label{eq:c}
    P[(X, \Pi)] = \frac{\exp\left(\frac{\varepsilon \cdot I(X; \Pi)}{2\Delta I}\right)}{\sum_{(X', \Pi') \in \Omega} \exp\left(\frac{\varepsilon \cdot I(X'; \Pi')}{2\Delta I}\right)}.
\end{equation}

\section{Генерация синтетических данных}

После построения структуры Байесовской сети выполняется генерация синтетических данных. Процесс включает два основных этапа:
\begin{enumerate}[label*=---]
	\item построение приватных условных распределений;
	\item выборка данных из построенных распределений.
\end{enumerate}

% \subsubsection{Построение условных распределений}

Для каждого узла в сети строится условное распределение $P(X_i | Parents(X_i))$ с добавлением шума Лапласа:

\begin{equation}
    \tilde{N}(x_i, pa_i) = N(x_i, pa_i) + \text{Lap}\left(\frac{k-1}{\varepsilon}\right),
\end{equation}
где:
\begin{enumerate}[label*=---]
	\item $N(x_i, pa_i)$ --- количество записей с значениями $x_i$ и $pa_i$;
	\item $k$ --- количество атрибутов в сети.
\end{enumerate}


% \subsubsection{Алгоритм генерации}

На начальном этапе генерации метод читает построенную БС и метаданные атрибутов.
Корневые атрибуты, которые не зависят от других переменных, сэмплируются из их безусловных распределений, а дочерние --- из условных.
Обход сети происходит в топологическом порядке.

\newpage
Схема алгоритма генерации синтетических данных предаставлена на рисунке~\ref{img:xz}:
\includeimage
{xz} % Имя файла без расширения (файл должен быть расположен в директории inc/img/)
{f} % Обтекание (с обтеканием)
{h!} % Положение рисунка (см. wrapfigure из пакета wrapfig)
{\textwidth} % Ширина рисунка
{Схема алгоритма генерации синтетических данных} % Подпись рисунка


\section{Особенности реализации}

\subsubsection{Обработка категориальных атрибутов}

Категориальные атрибуты кодируются в целочисленные значения для эффективной работы с Байесовской сетью:
\begin{equation}
    encode(x_i) = \begin{cases}
        value\_to\_index[x_i], & \text{если атрибут категориальный}, \\
        bin\_index(x_i), & \text{если атрибут непрерывный}.
    \end{cases}
\end{equation}

\subsubsection{Обработка пропущенных значений}

Пропущенные значения обрабатываются как отдельная категория с индексом $|domain| + 1$. Вероятность генерации пропущенного значения определяется по формуле:
\begin{equation}
    P(\text{missing}) = \frac{count(\text{missing values})}{total\_count}.
\end{equation}

\subsubsection{Постобработка результатов}

После генерации закодированных значений выполняется их преобразование в исходный формат:
\begin{enumerate}[label*=---]
	\item декодирование категориальных значений;
	\item денормализация числовых значений;
	\item форматирование временных меток;
	\item генерация строковых значений заданной длины.
\end{enumerate}

% \section{Анализ сложности}

% Временная сложность разработанного алгоритма:
% \begin{itemize}[label*=---]
%	\item Предобработка: $O(n \cdot d)$, где $n$ --- количество записей, $d$ --- количество атрибутов;
%	\item Построение БС: $O(d^3 \cdot k \cdot n)$, где $k$ --- максимальное количество родителей;
%	\item Генерация данных: $O(m \cdot d)$, где $m$ --- количество генерируемых записей.
% \end{itemize}

% Общая сложность: $O(d^3 \cdot k \cdot n + m \cdot d)$.

% Пространственная сложность: $O(n \cdot d + d^2)$ для хранения исходных данных и структуры сети.